\documentclass[../../../include/open-logic-section]{subfiles}

\begin{document}

\olfileid{sth}{choice}{woproblem}
\olsection{良序问题}

显然，我们是否接受良序公理关系重大。
但关于这一原则的讨论，往往聚焦于一个等价的原则，即选择公理。
因此，我们现在将把注意力转向后者（并证明其等价性）。

在 \citeyear{Cantor1883}，Cantor（康托尔）表达了他对良序公理的支持，称其为“一个在我看来是基本的、蕴含丰富的、且因其普遍有效性而尤为引人注目的思维法则”（转引自 \citeauthor{Potter2004} \citeyear[p.~243]{Potter2004}）。但 Cantor 最终确信，这个“公理”需要证明。数学界也是如此。

Zermelo（策梅洛）在 \citeyear{Zermelo1904} “解决”了这个问题。
为了解释他的解决方案，我们需要一些定义。

\begin{defn}
一个函数 $f$ 是一个\emph{选择函数}，当且仅当对所有 $x \in \dom{f}$，有 $f(x) \in x$。我们称 $f$ 是\emph{$A$ 的选择函数}，当且仅当 $f$ 是一个选择函数且满足 $\dom{f} = A \setminus \{\emptyset\}$。
\end{defn}

直观地说，对于 $A$ 中的每个（非空）集合 $x$，$A$ 的一个选择函数从 $x$ 中\emph{选出}一个特定的元素 $f(x)$。
那么，选择公理就是：

\begin{axiom}[选择]
	每个集合都有一个选择函数。
\end{axiom}

Zermelo 证明了选择公理能推出良序公理，反之亦然：

\begin{thm}[在 $\ZF$ 中]\ollabel{thmwochoice}
良序公理与选择公理等价。
\end{thm}

\begin{proof}
\emph{从左到右。}设 $A$ 是一个集合的集合。根据并集公理，$\bigcup A$ 存在，因此由良序公理，存在某个 $<$ 将 $\bigcup A$ 良序化。现在令 $f(x) = \text{the $<$-least
member of }x$。这就是 $A$ 的一个选择函数。

\emph{从右到左。}固定 $A$。由选择公理，存在 $\Pow{A} \setminus \{\emptyset\}$ 的一个选择函数 $f$。利用超限递归，定义一个函数：
\begin{align*}
	g(0) &= f(A)\\
	g(\alpha) &= 
		\begin{cases}
			\text{停止！} &\text{若}A = \funimage{g}{\alpha}\\
			f(A \setminus \funimage{g}{\alpha}) & \text{否则}\\	
		\end{cases}
\end{align*}
“停止！”的指示只是对原本可能更冗长的定义的一种简写。也就是说，当 $A = \funimage{g}{\alpha}$ 首次成立时，令对所有 $\delta \leq \alpha$ 有 $g(\delta) = A$。起初，我们只能确定这定义了一个\emph{项}（参见 \olref[sth][spine][recursion]{transrecursionschema} 后的说明）；但我们将证明这确实定义了一个函数。

由于 $f$ 是选择函数，对每个（当有定义时的）$\alpha$，我们有 $g(\alpha) = f(A \setminus \funimage{g}{\alpha}) \in A \setminus \funimage{g}{\alpha}$；即，$g(\alpha) \notin \funimage{g}{\alpha}$。
因此，如果 $g(\alpha) = g(\beta)$，那么 $g(\beta) \notin \funimage{g}{\alpha}$，即 $\beta \notin \alpha$，类似地有 $\alpha \notin \beta$。所以根据三歧性，$\alpha = \beta$。故 $g$ 是单射。

接着，注意到我们确实会停止！，即存在某个（最小的）序数 $\alpha$ 使得 $A = g[\alpha]$。因为若非如此；那么由于 $g$ 是单射，我们将对每个序数 $\alpha$ 有 $\cardless{\alpha}{\Pow{A} \setminus \{\emptyset\}}$，这与 \olref[hartogs]{HartogsLemma} 矛盾。因此也有 $\ran{g} = A$。

综合这些事实，$g$ 是从某个序数到 $A$ 的一个双射。现在 $g$ 可用来将 $A$ 良序化。
\end{proof}

所以，良序公理与选择公理同进退。
但问题依然存在：它们是成立，还是不成立呢？

\end{document}